\chapter{Doctoral Entrance Problems}
\label{ch:problems}

The problems below are of the type and difficulty of the linear algebra parts of
Algerian doctoral entrance examinations: a short statement, no hint, and a
complete proof expected in twenty to thirty minutes. Full solutions are in
\cref{ch:solutions}.
\newpage

\begin{problem}{(USTHB) — Épreuve N°01}{}
Soit l'ensemble $V = \{M \in \mathcal{M}_2(\mathbb{R}) \;/\; {}^tM = M\}$.
\begin{enumerate}
\item Montrer que $V$ est un sous-espace vectoriel de $\mathcal{M}_2(\mathbb{R})$.

\item Déterminer une base de $V$ ainsi que la dimension de $V$.

\end{enumerate}

On pose
$$
A=\begin{pmatrix}
1 & 1\\
0 & 2
\end{pmatrix}
$$
et pour toute matrice $M\in V$, on note
$$
f(M)=AM({}^tA).
$$

\begin{enumerate}
\setcounter{enumi}{2}

\item Montrer que $f$ est un endomorphisme de $V$.

\item Donner la matrice $B$ associée à $f$ dans la base de $V$ obtenue précédemment.

\item Déterminer le spectre de $f$.

\item L'endomorphisme $f$ est-il diagonalisable sur $\mathbb{R}$ ? Justifier.

\end{enumerate}

\source{\emph{\small Université des Sciences et de la Technologie Houari Boumediène (USTHB)}}
\end{problem}

\newpage

\begin{problem}{(USTHB) — Épreuve N°02}{antisymmetric-r3}
On munit $E=\mathbb{R}^3$ de sa base canonique $\mathcal{B}=\{e_1,e_2,e_3\}$ et de son produit scalaire usuel~:
pour tous $x=(x_1,x_2,x_3)$ et $y=(y_1,y_2,y_3)$ dans $E$,
$\langle x,y\rangle = x_1 y_1 + x_2 y_2 + x_3 y_3$.

Soit $u$ un endomorphisme non nul de $E$ tel que
$\langle u(x),x\rangle=0$ pour tout $x\in E$.
\begin{enumerate}
  \item Montrer que, pour tout $(x,y)\in E^{2}$,
        $\langle u(x),y\rangle = -\langle x,u(y)\rangle$.
        On pourra utiliser $z=x+y$ et calculer $\langle u(z),z\rangle$.
  \item \begin{enumerate}
          \item Démontrer que $0$ est la seule valeur propre réelle possible de $u$.
          \item Soit $P_u(X)$ le polynôme caractéristique de $u$. Quel est son degré~?
                Pourquoi a-t-il au moins une racine réelle~? En déduire que
                $P_u(X)=X\,R(X)$ où $R(X)$ est de degré~2 sans racines réelles.
          \item Montrer que $\ker(u)$ est de dimension~1. Quelle est la dimension de $\mathrm{Im}(u)$~?
        \end{enumerate}
  \item \begin{enumerate}
          \item Démontrer que si $x\in\ker(u)$ et $y\in\mathrm{Im}(u)$ alors
                $\langle x,y\rangle=0$.
          \item En déduire que $E$ est somme directe de $\ker(u)$ et $\mathrm{Im}(u)$.
        \end{enumerate}
\end{enumerate}

\source{\emph{\small Université des Sciences et de la Technologie Houari Boumediène (USTHB), 2026, Épreuve~n\textsuperscript{o}~02.}}
\end{problem}
